% ---------------------------------------------------------------------------
%  Lecture notes of Thermodynamics (ENGI 1111) -- shared LaTeX preamble
% ---------------------------------------------------------------------------

% The text font is set in _pdf.yml (fontfamily: mathptmx) so that Pandoc's own
% template loads it; do not load a second text font here.

\usepackage{amsmath,amssymb}
\usepackage{xcolor}
\usepackage{tikz}
\usepackage{fancyhdr}
\usepackage[most]{tcolorbox}
\usepackage{enumitem}
\usepackage{microtype}
\usepackage{wrapfig}
\usepackage{graphicx}

% KOMA sets headings in sans by default; the original Word notes are serif
\setkomafont{disposition}{\rmfamily\bfseries}
\setkomafont{title}{\rmfamily\bfseries}
\setkomafont{subtitle}{\rmfamily\mdseries}

% --- Durham palette --------------------------------------------------------
\definecolor{duPurple}{HTML}{68246D}
\definecolor{duBlue}{HTML}{002A41}
\definecolor{duRed}{HTML}{BE1E2D}
\definecolor{duGrey}{HTML}{6B6B6B}
\definecolor{duGreen}{HTML}{1B7A57}
\definecolor{duPaper}{HTML}{F4F0F5}

% --- running heads ---------------------------------------------------------
\newcommand{\gapvariant}{}          % overwritten per build by filters/gaps.lua

\pagestyle{fancy}
\fancyhf{}
\setlength{\headheight}{14pt}
\fancyhead[L]{\footnotesize\color{duGrey} Lecture notes of Thermodynamics (ENGI 1111)}
\fancyhead[R]{\footnotesize\color{duGrey} Instructor: Majid Bastankhah}
\fancyfoot[L]{\footnotesize\color{duGrey}\gapvariant}
\fancyfoot[R]{\footnotesize\color{duGrey}\thepage}
\renewcommand{\headrulewidth}{0.4pt}
\renewcommand{\headrule}{\hbox to\headwidth{\color{duPurple}\leaders\hrule height \headrulewidth\hfill}}

% --- gaps ------------------------------------------------------------------
% student copy: a ruled blank the students fill in with you
\newcommand{\gapblank}[1]{%
  \par\vspace{0.6em}%
  \noindent
  \begin{tikzpicture}
    \draw[duPurple!25, dashed, rounded corners=2pt]
      (0,0) rectangle (\linewidth,#1);
  \end{tikzpicture}%
  \par\vspace{0.6em}%
}

% completed copy: what was written in class, flagged with a purple rule
\newtcolorbox{gapfilled}{
  enhanced, breakable, colback=duPurple!3, colframe=duPurple!55,
  boxrule=0pt, leftrule=2pt, arc=0pt, outer arc=0pt,
  left=8pt, right=6pt, top=6pt, bottom=6pt,
  before skip=8pt, after skip=8pt
}

% lecturer copy: the same content, small and grey, so nothing is improvised
\newtcolorbox{gaplecturer}{
  enhanced, breakable, colback=white, colframe=duGrey!45,
  boxrule=0pt, leftrule=1pt, arc=0pt,
  borderline west={1pt}{0pt}{duGrey!45, dashed},
  left=8pt, right=6pt, top=4pt, bottom=4pt,
  before skip=6pt, after skip=6pt,
  fontupper=\small\color{duGrey!160}
}
\newcommand{\gaphintlabel}{{\bfseries\color{duPurple}Board:}\ \itshape}

% --- content boxes ---------------------------------------------------------
\newtcolorbox{tnote}{
  enhanced, breakable, colback=duBlue!4, colframe=duBlue!30,
  boxrule=0.4pt, arc=2pt, left=7pt, right=7pt, top=5pt, bottom=5pt,
  before skip=7pt, after skip=7pt
}
\newtcolorbox{twarning}{
  enhanced, breakable, colback=duRed!4, colframe=duRed!45,
  boxrule=0.4pt, arc=2pt, left=7pt, right=7pt, top=5pt, bottom=5pt,
  before skip=7pt, after skip=7pt
}
\newtcolorbox{tkey}{
  enhanced, breakable, colback=duPurple!7, colframe=duPurple,
  boxrule=0.8pt, arc=2pt, left=7pt, right=7pt, top=6pt, bottom=6pt,
  before skip=8pt, after skip=8pt
}
\newtcolorbox{texample}{
  enhanced, breakable, colback=white, colframe=duPurple!70,
  boxrule=0.6pt, arc=2pt, left=7pt, right=7pt, top=6pt, bottom=6pt,
  before skip=9pt, after skip=9pt
}
\newtcolorbox{tactivity}{
  enhanced, breakable, colback=duPurple!10, colframe=duPurple,
  boxrule=0.8pt, arc=2pt, left=7pt, right=7pt, top=6pt, bottom=6pt,
  before skip=9pt, after skip=9pt
}
\newtcolorbox{tsolution}{
  enhanced, breakable, colback=duGreen!5, colframe=duGreen!55,
  boxrule=0.4pt, arc=2pt, left=7pt, right=7pt, top=6pt, bottom=6pt,
  before skip=7pt, after skip=10pt
}
\newtcolorbox{tobjectives}{
  enhanced, breakable, colback=duPaper, colframe=duPurple!40,
  boxrule=0.4pt, arc=2pt, left=7pt, right=7pt, top=6pt, bottom=6pt,
  before skip=6pt, after skip=10pt
}

% labels used by filters/boxes.lua
\newcommand{\boxlabel}[1]{{\bfseries\color{duPurple}#1}\quad}

% --- misc ------------------------------------------------------------------
\setlist[itemize]{itemsep=2pt, topsep=3pt}
\setlist[enumerate]{itemsep=2pt, topsep=3pt}

% NOTE: no custom math macros are defined here on purpose. The same maths has
% to render in the HTML build too (MathJax), so everything in the .qmd files
% uses plain \mathrm{d} and P_\text{ext} rather than shorthands.
